% Advanced Quantum Technologies working front matter
% Source manuscript: paper/manuscript/spectral_geometry_rewrite.tex
% Stage 1: title, keywords, abstract

\title{Spectral Orientation of Measurement-Accessible Quantum Tangent Geometry in Variational Quantum Circuits}

% AQT requires 3--7 keywords.
\textbf{Keywords:} variational quantum circuits; quantum information geometry; quantum measurements; readout design; Fisher information; quantum machine learning

\begin{abstract}
Variational quantum circuits are accessed through measurements, but learning pipelines often retain only a restricted set of observables or classical features. Readout dimension alone does not determine how much measurement-induced tangent information survives this restriction. A normalized covariance of tangent scores separates readout rank, spectral concentration, and readout orientation, while random rank-matched subspaces define a rank-only reference. Across five generic circuit families, one- and two-body diagonal readouts remain near this reference through 16 qubits although the tangent covariance becomes strongly anisotropic. In an equal-rank comparison for Haar-$U(4)$ circuits at 12 qubits, a cross-fitted aligned readout increases the mean directional gradient-energy proxy by a factor of 9.584 and the finite-shot signal-to-noise ratio by a factor of 3.111 relative to the physical one-body readout; a random rank-matched subspace remains near the reference. A half-filled $U(1)$-conserving family shows the complementary regime, where low-weight $Z$ readouts are already aligned with leading tangent directions. Readout orientation therefore controls measurement-accessible tangent information independently of readout rank and can materially alter finite-shot directional signal.
\end{abstract}
